
pour tout @@M0@@, et on a
@@M1@@
Alors l'assertion (1.1.3) pour \emph{tous} les membres de (1.1.5) découle immédiatement de l'hypothèse et de l'équation
@@M2@@

On sait que @@M3@@ est un foncteur covariant sur la catégorie des schémas localement noethériens par rapport aux morphismes propres. La proposition suivante affirme que notre filtration est fonctorielle :

\textbf{Proposition 1.1.8.} \uline{Soit} @@M4@@ \uline{un morphisme propre de schémas localement noethériens et soit}
@@M5@@
\uline{le morphisme induit par} @@M6@@\uline{. Alors}
@@M7@@
\uline{pour tout entier} @@M8@@.

\emph{Démonstration :} Si @@M9@@ est un @@M10@@-module cohérent, alors on a par définition (IV 2.11) :
@@M11@@
Or, les @@M12@@-modules @@M13@@ sont concentrés sur @@M14@@ et d'après (EGA IV 5.4.1 (ii)) on a
@@M15@@
d'où l'assertion.

